\section{Schwarz inequalities for positive maps}

The following consequences of positivity and the Kadison--Schwarz inequality
concern normal, subnormal, and commuting-dominant operators.  They provide
matrix-order tools used independently of the Fundamental Theorem.

\subsection{Schwarz inequality for normal operators}

\begin{theorem}[CP Schwarz inequality for normal operators]\label{thm:ks_normal_cp}
    \lean{KadisonSchwarz.kadison_schwarz_normal_cp, KadisonSchwarz.kadison_schwarz_normal_cp_le}
    \leanok
    \uses{def:tp_kraus, def:kraus_adjoint_map}
    Let $E^*(X) = \sum_i K_i^\dagger X K_i$ be the adjoint Kraus map of a
    trace-preserving Kraus family, and let $A \in \MN{D}$ be normal.
    Then the Schwarz gap
    \begin{equation}\label{eq:schwarz_normal_cp_gap}
        E^*(A^\dagger A) - E^*(A^\dagger) E^*(A)
    \end{equation}
    is positive semidefinite. Equivalently,
    \begin{equation}\label{eq:schwarz_normal_cp_loewner}
        E^*(A^\dagger) E^*(A) \le E^*(A^\dagger A).
    \end{equation}
    This is the completely positive case of \cite[Proposition~5.1]{Wolf2012Quantum}.
\end{theorem}

\begin{proof}
    \leanok
    \uses{thm:kadison_schwarz_adjoint}
    The positivity statement is exactly Theorem~\ref{thm:kadison_schwarz_adjoint}.
    The second formulation~(\ref{eq:schwarz_normal_cp_loewner}) is the Loewner-order reformulation of the first~(\ref{eq:schwarz_normal_cp_gap}).
\end{proof}

\begin{theorem}[Schwarz inequality for normal operators]\label{thm:schwarz_normal_operator}
    \lean{KadisonSchwarz.schwarz_inequality_normal_operator}
    \leanok
    \uses{def:positive_map}
    Let $T : \MN{D} \to \MN{D}$ be a positive linear map with
    $T(\Id) \le \Id$.
    If $A \in \MN{D}$ is normal, then
    \[
        T(A^\dagger A) - T(A^\dagger) T(A) \ge 0.
    \]
    This is \cite[Proposition~5.1]{Wolf2012Quantum}.
\end{theorem}

\begin{proof}
    \leanok
    Since $A$ is normal, diagonalize $A = U D U^\dagger$ with $D$ diagonal
    and $U$ unitary.
    Express $A^\dagger A = U |D|^2 U^\dagger$ as a sum
    $\sum_{i,j} \bar{d}_i d_j P_i P_j$ where $P_i = U e_i e_i^\dagger U^\dagger$.
    Each $P_i$ is rank-one PSD, and the images $T(P_i)$ pairwise commute
    because they are images of mutually orthogonal projectors under a
    positive map.
    The Schwarz gap then reduces to a block quadratic form in the images,
    which is shown non-negative by a simultaneous-diagonalization argument:
    positivity of $T$ ensures the scalar PSD matrix
    $(\langle v_k, T(P_i P_j) v_k \rangle)_{i,j}$ dominates the product
    $(\langle v_k, T(P_i) v_k \rangle \langle v_k, T(P_j) v_k \rangle)_{i,j}$,
    and combining over a common eigenbasis of the $T(P_i)$ yields the result.
\end{proof}

\begin{remark}
    The proof follows the positivity-on-abelian argument of
    \cite[Proposition~1.6]{Wolf2012Quantum}: a positive map restricted to the
    commutative $*$-subalgebra generated by the spectral projectors of $A$
    satisfies the Schwarz inequality directly, without invoking the
    Arveson extension theorem.
\end{remark}

\subsection{Schwarz inequality for subnormal and commuting-dominant operators}

\begin{theorem}[Schwarz inequality for subnormal operators]\label{thm:schwarz_subnormal}
    \lean{KadisonSchwarz.schwarz_inequality_subnormal_operator}
    \leanok
    \uses{def:positive_map}
    Let $T : \MN{D} \to \MN{D}$ be a positive linear map with
    $T(\Id) \le \Id$.
    If $A \in \MN{D}$ is subnormal (i.e.\ there exists a normal operator
    on a larger space whose compression to $\C^D$ is $A$), then
    \[
        T(A^\dagger A) - T(A^\dagger) T(A) \ge 0.
    \]
    This is \cite[Theorem~5.5]{Wolf2012Quantum}.
\end{theorem}

\begin{proof}
    \leanok
    \uses{thm:schwarz_normal_operator}
    Embed $A$ as the northwest corner of a normal operator $N$ on
    $\C^D \oplus \C^E$.
    Extend $T$ to a positive subunital map $\tilde{T}$ on the larger
    space (acting as $T$ on the northwest block, zero elsewhere).
    Apply Theorem~\ref{thm:schwarz_normal_operator} to $\tilde{T}$ and
    $N$, then compress the resulting inequality back to $\C^D$.
\end{proof}

\begin{theorem}[Schwarz inequality for commuting-dominant operators]\label{thm:schwarz_commuting_dominant}
    \lean{KadisonSchwarz.schwarz_inequality_commuting_dominant_operator}
    \leanok
    \uses{def:positive_map}
    Let $T : \MN{D} \to \MN{D}$ be a positive linear map with
    $T(\Id) \le \Id$.
    If $A \in \MN{D}$ admits a positive semidefinite \emph{dominant}
    $B \ge 0$ with $A^\dagger A \le B$ and $B$ commuting with $A$
    ($BA = AB$),
    then
    \[
        T(A^\dagger) T(A) \le T(B)
        \quad\text{and}\quad
        T(A) T(A^\dagger) \le T(B).
    \]
    This is \cite[Theorem~5.6]{Wolf2012Quantum}.
\end{theorem}

\begin{proof}
    \leanok
    \uses{thm:schwarz_normal_operator}
    First take $B > 0$. Since $B$ commutes with $A$ and $A^\dagger A \le B$,
    also $A A^\dagger \le B$, so $D_R = (B - A^\dagger A)^{1/2}$ and
    $D_L = (B - A A^\dagger)^{1/2}$ are positive semidefinite. The block matrix
    $N = \left(\begin{smallmatrix} A & D_L \\ D_R & -A^\dagger \end{smallmatrix}\right)$
    is normal, with $N^\dagger N = N N^\dagger = B \oplus B$. As in the
    subnormal case, extend $T$ to a positive subunital map on $\MN{2D}$, apply
    the normal-operator Schwarz inequality
    (Theorem~\ref{thm:schwarz_normal_operator}) to that extension and $N$, and
    compress back to $\C^D$. Since $N$ is normal, both
    $T(A^\dagger) T(A) \le T(B)$ and $T(A) T(A^\dagger) \le T(B)$ follow. For
    general $B \ge 0$, apply this to $B_\varepsilon = B + \varepsilon \Id > 0$
    and let $\varepsilon \to 0$.
\end{proof}

\begin{theorem}[CP variant: Kadison--Schwarz for commuting-dominant inputs]\label{thm:ks_commuting_dominant_cp}
    \lean{KadisonSchwarz.kadison_schwarz_commuting_dominant_cp}
    \leanok
    \uses{def:tp_kraus, def:kraus_adjoint_map}
    Let $E^*$ be the adjoint of a trace-preserving Kraus map.
    If $A \in \MN{D}$ admits a positive semidefinite dominant $B$
    commuting with $A$ and satisfying $A^\dagger A \le B$,
    then $E^*(A^\dagger) E^*(A) \le E^*(B)$ and
    $E^*(A) E^*(A^\dagger) \le E^*(B)$.
\end{theorem}

\begin{proof}
    \leanok
    \uses{thm:schwarz_commuting_dominant}
    Apply Theorem~\ref{thm:schwarz_commuting_dominant} to the adjoint
    Kraus map (which is unital under the trace-preservation hypothesis).
\end{proof}

\section{Diagonal Jensen inequality}

\begin{lemma}[Diagonal Jensen inequality]\label{lem:diagonal_jensen_of_convexOn}
    \lean{Matrix.diagonal_jensen_of_convexOn}
    \leanok
    Let $f:\R\to\R$ be convex on $[0,\infty)$, let $A$ be a positive
    semidefinite matrix on a finite-dimensional space indexed by $n$, and let $v\in\C^n$
    satisfy $\langle v,v\rangle=1$. Then
    \[
        f\!(\Re\langle v, A v\rangle)
        \le
        \Re\langle v, f(A)\,v\rangle,
    \]
    where $f(A)$ is defined by the Hermitian continuous functional calculus.
\end{lemma}

\begin{proof}\leanok
    Write $A=U\,\diag(\mu)\,U^{\ast}$ by the spectral
    theorem and set $w=U^{\ast}v$. Since $U$ is unitary,
    $p_i:=\lvert w_i\rvert^2$ satisfies $\sum_i p_i=\lVert v\rVert^2=1$,
    so $(p_i)$ is a probability distribution on $n$. The eigenvalues
    $\mu_i$ lie in $[0,\infty)$ because $A$ is positive semidefinite.
    A direct computation yields
    $\Re\langle v,Av\rangle=\sum_i p_i\mu_i$ and
    $\Re\langle v,f(A)v\rangle=\sum_i p_i f(\mu_i)$, so the scalar Jensen
    inequality applied to the weights $(p_i)$ and points $(\mu_i)$ gives
    the conclusion. See \cite[Chapter~V]{Bhatia1997Matrix}.
\end{proof}

\section{Trace concavity and convexity of matrix powers}

\begin{theorem}[Trace concavity of real powers]\label{thm:trace_rpow_concave}
    \lean{TNLean.trace_rpow_concave, trace_rpow_concave}
    \leanok
    \uses{lem:diagonal_jensen_of_convexOn}
    For $p\in[0,1]$, PSD matrices $A_1,A_2$, and $t\in[0,1]$:
    \[
        t\,\Re\tr(A_1^p)+(1-t) \Re\tr(A_2^p)
        \le
        \Re\tr\!((tA_1+(1-t)A_2)^p).
    \]
    Follows from operator concavity of $x\mapsto x^p$ composed with trace
    monotonicity on the Loewner order.
\end{theorem}

\begin{proof}\leanok
    Let $A=tA_1+(1-t)A_2$ and let $\{\psi_j\}$ with eigenvalues
    $\mu_j\ge 0$ be an orthonormal eigenbasis of $A$. A direct
    computation gives $\Re\tr(A^p)=\sum_j \mu_j^p$ and
    $\mu_j=t\,a_j+(1-t)\,b_j$ where
    $a_j=\Re\langle\psi_j,A_1\psi_j\rangle$,
    $b_j=\Re\langle\psi_j,A_2\psi_j\rangle$.
    Scalar concavity of $x\mapsto x^p$ on $[0,\infty)$ yields
    $\mu_j^p\ge t\,a_j^p+(1-t)\,b_j^p$, and the diagonal Jensen
    inequality (Lemma~\ref{lem:diagonal_jensen_of_convexOn},
    concave form) gives
    $a_j^p\le\Re\langle\psi_j,A_1^p\psi_j\rangle$ (and similarly for
    $b_j$). Summing over $j$ and using
    $\sum_j\Re\langle\psi_j,M\psi_j\rangle=\Re\tr M$ yields the
    conclusion. See \cite[Chapter~V]{Bhatia1997Matrix}.
\end{proof}

\begin{theorem}[Trace convexity of real powers]\label{thm:trace_rpow_convex}
    \lean{TNLean.trace_rpow_convex, trace_rpow_convex}
    \leanok
    \uses{lem:diagonal_jensen_of_convexOn, thm:trace_rpow_concave}
    For $p\in[1,2]$, PSD matrices $A_1,A_2$, and $t\in[0,1]$:
    \[
        \Re\tr\!((tA_1+(1-t)A_2)^p)
        \le
        t\,\Re\tr(A_1^p)+(1-t) \Re\tr(A_2^p).
    \]
    Follows from operator convexity of $x\mapsto x^p$ composed with trace
    monotonicity on the Loewner order.
\end{theorem}

\begin{proof}\leanok
    Identical eigenbasis reduction as in
    Theorem~\ref{thm:trace_rpow_concave}, with scalar convexity
    (rather than concavity) of $x\mapsto x^p$ on $[0,\infty)$ for
    $p\ge 1$ and the convex form of the diagonal Jensen inequality.
\end{proof}

\section{Operator convexity of real powers}

\begin{theorem}[Operator convexity of real powers between one and two]%
    \label{thm:operator_convexOn_rpow}
    \lean{CFC.convexOn_rpow}
    \leanok
    For elements $a,b$ of a unital C${}^\ast$-algebra with $0\le a,b$, an
    exponent $p\in[1,2]$, and $t\in[0,1]$, the map $a\mapsto a^p$ is convex for
    the Loewner order on the positive cone:
    \[
        (t a+(1-t) b)^p \le t\,a^p+(1-t)\,b^p.
    \]
    This is the convex counterpart of the operator concavity of $x\mapsto x^p$
    on $[0,1]$.
\end{theorem}

\begin{proof}\leanok
    For $p\in(1,2)$ one uses the integral representation
    $a^p=\int_0^\infty f_{p,s}(a)\,d\mu(s)$ of Carlen
    \cite[Lemma~2.8]{carlen2010}, whose integrand
    $f_{p,s}(x)=s^{p-1}\bigl(s^{-1}x+s(s+x)^{-1}-1\bigr)$ is operator convex:
    it decomposes into a linear term, a nonnegative multiple of the operator
    convex resolvent $x\mapsto(s\cdot 1+x)^{-1}$, and a constant.
    Since $f_{p,s}(ta+(1-t)b)\le t\,f_{p,s}(a)+(1-t)\,f_{p,s}(b)$ for every
    $s>0$, integrating this pointwise Loewner inequality against $\mu$ gives
    \[
        (ta+(1-t)b)^p
        =\int_0^\infty f_{p,s}(ta+(1-t)b)\,d\mu(s)
        \le t\int_0^\infty f_{p,s}(a)\,d\mu(s)
            +(1-t)\int_0^\infty f_{p,s}(b)\,d\mu(s)
        = t\,a^p+(1-t)\,b^p.
    \]
    The endpoint $p=1$ is the identity, and $p=2$ is $a\mapsto a^2$, operator
    convex because
    $t\,a^2+(1-t)b^2-(ta+(1-t)b)^2=t(1-t)(a-b)^2\ge 0$.
    See \cite[Chapter~V]{Bhatia1997Matrix}.
\end{proof}

\section{Operator Jensen inequality for positive maps}

\begin{definition}[Finite-POVM dilation construction]\label{def:operator_jensen_povm_aux}
    \lean{TNLean.OperatorJensen.povmIsometry, TNLean.OperatorJensen.povmDiagonal}
    \leanok
    For a finite family of matrices $\{C_i\}_{i \in \iota}$ and a defect block
    $S$, one forms an isometric dilation whose rows are the adjoints of the
    matrices $C_i$ together with the adjoint of $S$. One also forms the scalar
    block-diagonal matrix whose $\iota$-blocks carry prescribed weights and
    whose defect block carries a single scalar.
\end{definition}

\begin{lemma}[Compression auxiliary for Jensen]\label{lem:operator_jensen_inverse_compression}
    \lean{Matrix.PosDef.inverse_compression_le}
    \leanok
    Let $Y$ be positive definite and let $W$ be an isometry. Then
    \begin{equation}\label{eq:opconv_inverse_compression_le}
        (W^\ast Y W)^{-1} \le W^\ast Y^{-1} W.
    \end{equation}
\end{lemma}

\begin{proof}
    \leanok
    Apply the Schur-complement criterion to the block matrix
    $\begin{psmallmatrix} Y^{-1} & W \\ W^\ast & W^\ast Y W \end{psmallmatrix}$.
    The lower-right Schur complement is zero, hence positive semidefinite, and
    compressing the resulting upper-left Schur complement by $W$ gives the
    inequality~(\ref{eq:opconv_inverse_compression_le}).
\end{proof}

\begin{lemma}[Finite-POVM compression identities]\label{lem:operator_jensen_povm_compression}
    \lean{TNLean.OperatorJensen.povmIsometry_star_mul,
        TNLean.OperatorJensen.povmIsometry_compress_diagonal,
        TNLean.OperatorJensen.povm_sum_add_defect,
        TNLean.OperatorJensen.povmDiagonal_posDef,
        TNLean.OperatorJensen.povmDiagonal_inv,
        TNLean.OperatorJensen.povmIsometry_compress_diagonal_inv}
    \leanok
    If the defect block satisfies
    \begin{equation}\label{eq:opconv_defect_block_relation}
        S S^\ast = \Id - \sum_{i \in \iota} C_i C_i^\ast,
    \end{equation}
    then the dilation is an isometry, compressing the weighted scalar
    block-diagonal matrix gives
    \[
        \sum_{i \in \iota} w_i\, C_i C_i^\ast + t\, S S^\ast,
    \]
    the defect relation rewrites as
    \[
        \sum_{i \in \iota} C_i C_i^\ast + S S^\ast = \Id,
    \]
    and strict positivity of all weights together with $t>0$ implies positive
    definiteness of the scalar block-diagonal matrix. If all weights are nonzero
    and the defect scalar is nonzero, the inverse diagonal has reciprocal
    weights, and compressing this inverse gives
    \[
        \sum_{i \in \iota} w_i^{-1}\, C_i C_i^\ast + t^{-1}\, S S^\ast.
    \]
\end{lemma}

\begin{proof}
    \leanok
    Expand the block matrix multiplication entrywise. The defect
    identity~(\ref{eq:opconv_defect_block_relation}) gives the isometry relation
    after summing the $\iota$-blocks and the defect block, and the weighted
    compression identity follows from the same expansion with
    the scalar diagonal inserted. Positive definiteness uses strict positivity
    of every diagonal entry, while the inverse formula uses the nonvanishing of
    every diagonal entry.
\end{proof}

\begin{lemma}[Finite-POVM resolvent inequality]\label{lem:operator_jensen_povm_resolvent}
    \lean{TNLean.OperatorJensen.povm_resolvent_inv_le}
    \leanok
    \uses{lem:operator_jensen_povm_compression, lem:operator_jensen_inverse_compression}
    Let $w_i\ge 0$ and $t>0$. If the defect block satisfies
    \[
        S S^\ast = \Id - \sum_{i \in \iota} C_i C_i^\ast,
    \]
    then
    \begin{equation}\label{eq:opconv_povm_resolvent_bound}
        \left(\sum_{i \in \iota} w_i\, C_i C_i^\ast + t\,\Id\right)^{-1}
        \le
        \sum_{i \in \iota} (w_i+t)^{-1}\, C_i C_i^\ast + t^{-1}\, S S^\ast.
    \end{equation}
\end{lemma}

\begin{proof}
    \leanok
    \uses{lem:operator_jensen_povm_compression, lem:operator_jensen_inverse_compression}
    Apply the compression inequality to the positive definite scalar
    block-diagonal matrix with entries $w_i+t$ and defect entry $t$. The
    compression identities rewrite the compressed matrix as
    $\sum_i w_i C_i C_i^\ast + t\Id$ and its inverse compression as the upper
    bound~(\ref{eq:opconv_povm_resolvent_bound}).
\end{proof}

\begin{lemma}[Spectral finite-POVM resolution]\label{lem:operator_jensen_spectral_resolution}
    \lean{TNLean.OperatorJensen.posSemidef_spectral_resolution}
    \leanok
    Let $A\ge 0$ be a finite matrix. There are positive semidefinite matrices
    $P_i$ and nonnegative real numbers $\lambda_i$ such that
    \[
        \sum_i P_i=\Id,
        \qquad
        A=\sum_i \lambda_i P_i.
    \]
\end{lemma}

\begin{proof}
    \leanok
    Diagonalize the Hermitian matrix $A$ by a unitary matrix. The matrices
    $P_i$ are the conjugates of the rank-one diagonal matrix units, their sum is
    the identity, and positive semidefiniteness of $A$ gives $\lambda_i\ge0$.
\end{proof}

\begin{lemma}[Positive-map resolvent inequality]\label{lem:operator_jensen_positive_map_resolvent}
    \lean{TNLean.OperatorJensen.positiveMap_resolvent_inv_le}
    \leanok
    \uses{lem:operator_jensen_povm_resolvent,
        lem:operator_jensen_spectral_resolution}
    Let $T$ be a positive subunital map, let $A\ge0$, and let $t>0$. Then
    \[
        (TA+t\Id)^{-1}
        \le
        T((A+t\Id)^{-1})+t^{-1}(\Id-T\Id).
    \]
\end{lemma}

\begin{proof}
    \leanok
    \uses{lem:operator_jensen_povm_resolvent,
        lem:operator_jensen_spectral_resolution}
    Use the spectral resolution $A=\sum_i\lambda_i P_i$. The matrices
    $T(P_i)$ are positive semidefinite and sum to $T\Id\le\Id$. Write each
    $T(P_i)$ as $C_iC_i^\ast$ using the positive square root, and write the
    defect $\Id-\sum_iT(P_i)$ as $SS^\ast$. The finite-POVM resolvent
    inequality gives the desired estimate after substituting the spectral
    formula for $(A+t\Id)^{-1}$.
\end{proof}

\begin{lemma}[Positive-map real-power integrand inequality]
    \label{lem:operator_jensen_positive_map_rpow_integrand}
    \lean{TNLean.OperatorJensen.positiveMap_rpowIntegrand₀₁_jensen}
    \leanok
    \uses{lem:operator_jensen_positive_map_resolvent}
    Let $T$ be a positive subunital map, let $A\ge0$, let
    $p\in(0,1)$, and let $t>0$. Then
    \[
        T\!\left(f_{p,t}(A)\right)
        \le
        f_{p,t}(TA),
        \qquad
        f_{p,t}(x)=t^p\left(t^{-1}-(t+x)^{-1}\right).
    \]
\end{lemma}

\begin{proof}
    \leanok
    \uses{lem:operator_jensen_positive_map_resolvent}
    Rewrite the continuous functional calculus for $f_{p,t}$ in resolvent
    form,
    \[
        f_{p,t}(X)=t^{p-1}\Id-t^p(t\Id+X)^{-1}.
    \]
    After applying the linear map $T$, the difference
    $f_{p,t}(TA)-T(f_{p,t}(A))$ is exactly $t^p$ times the positive matrix
    supplied by the positive-map resolvent inequality.
\end{proof}

\begin{lemma}[Positive-map convex real-power integrand inequality]
    \label{lem:operator_jensen_positive_map_rpow_integrand_convex}
    \lean{TNLean.OperatorJensen.positiveMap_rpowIntegrand₁₂_jensen}
    \leanok
    \uses{lem:operator_jensen_positive_map_resolvent}
    Let $T$ be a positive subunital map, let $A\ge0$, let
    $p\in(1,2)$, and let $t>0$. Then
    \[
        g_{p,t}(TA)
        \le
        T\!\left(g_{p,t}(A)\right),
        \qquad
        g_{p,t}(x)=t^{p-2}x+t^p(t+x)^{-1}-t^{p-1}.
    \]
    The inequality is reversed relative to the concave integrand, since
    $g_{p,t}$ is operator convex.
\end{lemma}

\begin{proof}
    \leanok
    \uses{lem:operator_jensen_positive_map_resolvent}
    Rewrite the continuous functional calculus for $g_{p,t}$ in resolvent form,
    \[
        g_{p,t}(X)=t^{p-2}X+t^p(t\Id+X)^{-1}-t^{p-1}\Id.
    \]
    Under the linear map $T$ the leading $t^{p-2}$ terms cancel and the constant
    $t^{p-1}$ terms combine, so the difference
    $T(g_{p,t}(A))-g_{p,t}(TA)$ is exactly $t^p$ times the positive matrix
    supplied by the positive-map resolvent inequality (using
    $t^p\,t^{-1}=t^{p-1}$).
\end{proof}

\begin{lemma}[Positive matrix-valued integrals]\label{lem:operator_jensen_integral_order}
    \lean{TNLean.OperatorJensen.integral_nonneg_matrix_of_ae}
    \leanok
    Let $f$ be a function from a measure space to a finite matrix algebra. If
    $f(x)\ge0$ for almost every $x$, then
    $\int f(x)\,d\mu(x)\ge0$.
\end{lemma}

\begin{proof}
    \leanok
    The Loewner order has closed positive cone in finite dimension. Hence the
    ordered Bochner integral preserves almost-everywhere nonnegativity. The
    corresponding almost-everywhere monotonicity statement is the general
    ordered Bochner integral theorem.
\end{proof}

\begin{theorem}[Jensen for concave real powers]\label{thm:posMap_rpow_concave_jensen}
    \lean{posMap_rpow_concave_jensen}
    \leanok
    \uses{def:positive_map, lem:operator_jensen_positive_map_rpow_integrand,
        lem:operator_jensen_integral_order}
    For a positive subunital map $T$, a positive semidefinite matrix~$A$,
    and $p\in[0,1]$:
    \[
        T(A^p) \le (TA)^p.
    \]
    This is \cite[Theorem~5.13]{Wolf2012Quantum} for concave~$f(x)=x^p$.
\end{theorem}

\begin{proof}\leanok
    For $p=0$ the assertion is the subunitality condition $T(\Id)\le \Id$.
    For $p=1$ it is an equality.  For $0<p<1$, use Carlen's integral
    representation of $A^p$ by the functions
    $f_{p,t}(x)=t^p(t^{-1}-(t+x)^{-1})$.  The pointwise inequality
    $T(f_{p,t}(A))\le f_{p,t}(TA)$ is Lemma~\ref{lem:operator_jensen_positive_map_rpow_integrand}.
    Integrating this inequality with respect to the representing measure gives
    the result by Lemma~\ref{lem:operator_jensen_integral_order} and linearity
    of the Bochner integral under $T$.
\end{proof}

\begin{theorem}[Map form of Jensen for concave real powers]\label{thm:rpow_concave_jensen}
    \lean{IsPositiveMap.rpow_concave_jensen}
    \leanok
    \uses{thm:posMap_rpow_concave_jensen}
    For a positive subunital map $T$, a positive semidefinite matrix~$A$,
    and $p\in[0,1]$, the inequality
    $T(A^p) \le (TA)^p$ holds.
\end{theorem}

\begin{proof}\leanok
    Direct application of Theorem~\ref{thm:posMap_rpow_concave_jensen}.
\end{proof}

\begin{theorem}[Jensen for convex real powers]\label{thm:posMap_rpow_convex_jensen}
    \lean{posMap_rpow_convex_jensen}
    \leanok
    \uses{def:positive_map,
        lem:operator_jensen_positive_map_rpow_integrand_convex,
        lem:operator_jensen_integral_order}
    For a positive subunital map $T$, a positive semidefinite matrix~$A$,
    and $p\in[1,2]$:
    \[
        (TA)^p \le T(A^p).
    \]
    This is \cite[Theorem~5.11]{Wolf2012Quantum} for the subunital
    operator-convex case $f(x)=x^p$.
\end{theorem}

\begin{proof}\leanok
    For $p=1$ the assertion is an equality.  For $1<p<2$, use Carlen's
    integral representation of $A^p$ by the convex integrands
    $g_{p,t}(x)=t^{p-2}x+t^p(t+x)^{-1}-t^{p-1}$.  The pointwise inequality
    $g_{p,t}(TA)\le T(g_{p,t}(A))$ is
    Lemma~\ref{lem:operator_jensen_positive_map_rpow_integrand_convex}.
    Integrating this inequality with respect to the representing measure and
    using Lemma~\ref{lem:operator_jensen_integral_order} together with linearity
    of the Bochner integral under $T$ gives the inequality for $1<p<2$.  The
    endpoint $p=2$ follows by taking the limit $q\to2^-$ in
    $(TA)^q\le T(A^q)$: since $M^q\to M^2$ in the operator norm as $q\to2^-$ for
    any positive semidefinite $M$, and the Loewner order is closed, this gives
    $(TA)^2\le T(A^2)$.
\end{proof}

\begin{theorem}[Map form of Jensen for convex real powers]\label{thm:rpow_convex_jensen}
    \lean{IsPositiveMap.rpow_convex_jensen}
    \leanok
    \uses{thm:posMap_rpow_convex_jensen}
    For a positive subunital map $T$, a positive semidefinite matrix~$A$,
    and $p\in[1,2]$, the inequality
    $(TA)^p \le T(A^p)$ holds.
\end{theorem}

\begin{proof}\leanok
    Direct application of the theorem.
\end{proof}

\begin{theorem}[Jensen for concave log]\label{thm:posMap_log_concave_jensen}
    \lean{posMap_log_concave_jensen}
    \leanok
    \uses{def:positive_map, thm:posMap_rpow_concave_jensen}
    For a positive \emph{unital} map $T$ and positive-definite~$A$:
    \[
        T(\log A) \le \log(TA).
    \]
    This is \cite[Theorem~5.13]{Wolf2012Quantum} for~$f=\log$.
    Requires unitality ($T(\Id)=\Id$), not merely subunitality.
\end{theorem}

\begin{proof}\leanok
    Apply Theorem~\ref{thm:posMap_rpow_concave_jensen} to $0<p<1$ and rewrite
    it as an inequality for $(A^p-\Id)/p$.  The continuous functional calculus
    gives $(A^p-\Id)/p\to \log A$ as $p\to0^+$, and the same limit holds for
    $T(A)$.  Since the Loewner order graph is closed in finite dimension, the
    limiting inequality is $T(\log A)\le \log(TA)$.
\end{proof}

\begin{theorem}[Map form of Jensen for the concave logarithm]\label{thm:log_concave_jensen}
    \lean{IsPositiveMap.log_concave_jensen}
    \leanok
    \uses{thm:posMap_log_concave_jensen}
    For a positive unital map $T$ and a positive-definite matrix~$A$,
    the inequality $T(\log A) \le \log(TA)$ holds.
\end{theorem}

\begin{proof}\leanok
    Direct application of Theorem~\ref{thm:posMap_log_concave_jensen}.
\end{proof}
